\pagenumbering{arabic}


\part{Precalculus}
\chapter{Functions}
\section{Function basics}

Before we start on the calculus topics, this part of the book will acquaint you with the mathematics you would need to understand in order to read this book. If you are familiar with the these topics, you may skip to the Part II on Calculus I. For the rest of you, the pace of this part will be very fast, it is recommended that you read up on the concepts you do not understand very well. We will start with the idea of a \textit{function}. By now, I am sure that many of you may have seen equations like these:
\begin{equation}
  x^2-5x+6=0
\end{equation}
\begin{equation}
  \sin(x)=0.5, 0\leq x \leq\frac{\pi}{2}
\end{equation}
\begin{equation}
  \label{line}
  y=mx+c
\end{equation}

\begin{figure}
  \centering
  \begin{tikzpicture}
    \begin{axis}[xlabel=\(x\), ylabel=\(y\), axis lines=left]
      \addplot[mark=none, color=blue]{x};
    \end{axis}
  \end{tikzpicture} 
  \caption{Graph of \(y=x\).}
\end{figure}

Out of the three equations, notice that \Cref{line} is special. There is not some fixed \(x\) values that you must plug in. Instead, the \(y\) value varies with the \(x\) value. Given a \(x\) value, the equation "produces" a \(y\) value. This brings us to the idea of a function. A function is simply something that takes in a value, and "produces" another value based on the input. You may have came across functions before, such as the trigonometric functions like \(\sin\) and \(\cos\). The value of \(\sin(x)\) varies with different values of \(x\). Of course, it is time for me to give you the definition of a function. 
\begin{definition}
  A \textbf{function} \(f\) is a mathematical object that assigns its inputs, or \textit{domain}, to its outputs, or \textit{range}. The notation \(f:A \to B\) means that the function has a domain \(A\) and range \(B\) \footnotemark. The notation \(f(a)\) where \(a \in A\) denotes the element of \(B\) to which \(a\) is assigned to.    
\end{definition}
\footnotetext{Note that A and B are sets.}
\begin{example}{Evaluating functions}{}
  Here are some examples of functions:
  \begin{equation}
    f(x)=x+1 \label{x+1}
  \end{equation}
  \begin{equation}
    \label{sin}
    g(x)=\sin(x)
  \end{equation}
  \begin{equation}
    \label{quadratic}
    h(x)=x^2-5x+6
  \end{equation}
  To evaluate the output of a function, you simply substitute the input to the expression that defines the function. Notice that for \Cref{x+1}, it is simply the line \(y=x+1\). Each value of \(f(x)\) for some \(x \in \mathbb{R}\) is just \(y\). Take \(x=1\) for example, \(f(1)=1+1=2\) and for the line \(y=x+1\), notice the \(y\) value here is also 2. The function assigns each \(x\) value to \(x+1\). Hence, in this case, \(f:\mathbb{R} \to \mathbb{R}\). For \Cref{sin} , the function takes any real \(x\) value and maps it to \(\sin(x)\). As a result, its domain would be \(\mathbb{R}\) and its range would be \([-1,1]\). Again, it is just like the graph \(y=\sin(x)\). 
  Lastly, the quadratic in \Cref{quadratic} is another real function(i.e. it has a real domain and range). Think about what happens when \(h(x)=0\), what value of \(x\) must you plug in to make \(h(x)=0\). Well, isn't it just the roots of the quadratic. Hence, the roots of a function are just the values \(x\) such that \(h(x)=0\). In this case, it would be 2 and 3. Thus, \(h(2)=h(3)=0\).
\end{example}

Actually, functions can be more than the usual \(f(x)=\text{something}\). Let's say you want to have a function that outputs "1" if the input is rational and "0" if it is not, you can define it as such: \footnote{This is actually the famous Dirichlet function that shatters certain intuitions on functions and will be discussed in later chapters.}
\begin{equation}
  f(x)=
  \begin{cases}
    1 & \textit{if}\ x \in \mathbb{Q} \\
    0 & \textit{if}\ x \notin \mathbb{Q}
  \end{cases}
\end{equation}


Here, \(f(\frac{1}{2})\) would be 1 but \(f(\sqrt{2})\) would be 0. (You would probably get a migraine from graphing it.) This is an example of having conditionals in a function. Lastly, functions may take more than one input and produce more than one output, examples include:
\begin{equation}
  f(x,y)=x^2+y^2 
  \label{circle}
\end{equation}
\begin{equation}
  \label{paramcircle}
  g(t)=(\cos(t), \sin(t))
\end{equation}

\Cref{circle}\footnote{As you will see in Part IV of the book, this is the \textit{parametric} equation of a circle.}
is an example of a \textit{multivariate} function, in other words, it takes in multiple inputs. If \(x=2\) and \(y=2\), then \(f(2,2)=8\). On the other hand, \Cref{paramcircle} does not have a single output. Instead, it has a tuple \footnote{A tuple is basically a set where order matters. You can treat it as if it is a coordinate.} as its output. Evaluating it at \(t=\pi\), notice that the output is \((-1, 0)\). What is interesting is that if you trace out its outputs on a plane for all \(t \in \mathbb{R}\), you will get \Cref{param}.
\begin{figure}[h]
  \centering
  \begin{tikzpicture}
    \begin{axis}[trig format=rad, axis equal, axis lines=left, ymax=1.5, ymin=-1.5]
      \addplot[color=blue, samples=100]({cos(x)}, {sin(x)});
    \end{axis}
  \end{tikzpicture}
  \tiny\caption{Plot of \Cref{paramcircle} for all \(t \in \mathbb{R}\)}.
  \label{param}
\end{figure}
\subsection{Simple pendulum}
In physics, we often deal with functions which enables us to describe many things whose behaviour is dependent on certain inputs. Take a pendulum for example. We can describe the angle of the pendulum to the normal line using the following formula\footnote{While not completely accurate, this formula is good enough for an example.}:
\begin{equation}
  \theta(t) = \theta_0\cos(\sqrt{\frac{g}{l}}t)
\end{equation}
Here, \(\theta(t)\) is the angle of the pendulum at time \(t\), \(\theta_0\) is the initial angle of the pendulum, \(T\) is the period of the pendulum, \(l\) is the length of the rod. (\(g\) is the gravitational field strength.) While this function may look sightly complicated, we can deduce a lot about it without computing much. Firstly, at \(t=0\), \(\theta(0)=\theta_0\), which makes sense because at the start, the pendulum is at its initial angle.
\begin{figure}[h]
  \centering
  \includegraphics[height=8cm, width=8cm]{Simple_gravity_pendulum.png}
  \caption{Figure of pendulum.}
  \label{pendulum}
\end{figure}


The periodic nature of the pendulum is also expressed in the cosine part of the function. Moreover, since \(\cos\) can only attain values between -1 and 1, the pendulum is guaranteed to be oscillating between \(-\theta_0\) and \(\theta_0\). This satisfies the conservation of energy because if this was not guaranteed and no energy is introduced into the system, energy must have been created or destroyed. 
\begin{challenge}
  This being said, in the real world, is energy lost to the surroundings? How could energy be lost to the surroundings? How would this affect the motion of the pendulum?
\end{challenge}


\begin{example}
  Let us calculate the period of the pendulum. We know the formula for the angle of the pendulum for any time \(t\). The period of a function \(\cos(ax)\) is simply \(\frac{2\pi}{a}\), hence the period of the pendulum is \(2\pi\sqrt{\frac{l}{g}}\). This results shows us that the period of the pendulum is \textbf{independent} of the initial angle. This property is called "isochronism" that Galileo discovered.
\end{example}


In the real world, we often add a \textit{damping factor} \(d(t)\) to \(\theta(t)\). A damped pendulum has the function describing its angle as \(\theta_d(t)=\theta(t)d(t)\). The damping factor is to reduce the amplitude of the pendulum as time passes, mainly because air resistance opposes the motion of the pendulum. 
\begin{challenge}
  From the following damping factor \(d(t)\), choose the most suitable one (you may want to use a graphing calculator to assist you):
  \begin{align*}
    e^t && e^{-t} && -t && a, a\in \mathbb{R}
  \end{align*}
\end{challenge}

\begin{problem}
  Given that \(x=2\) and \(y=\pi\), evaluate \(f(x)\) for:
  \begin{enumerate}[label=\alph*)]
    \item{\(f(x)=x^3+1\)}
    \item{\(f(x)=2^x\)}
    \item{\(f(x)= \begin{cases}
        x^4+1 & \textit{if} \ x<2 \\
        3^x & \textit{otherwise}
    \end{cases}\)}
  \item{\(f(x,y)=\frac{\sin(y)}{x}\)}
  \end{enumerate}

\end{problem}
\begin{problem}
  Given the following functions, write down \(f:A \to B\) where \(A\) is the domain and \(B\) is the range(do not consider complex inputs and outputs, we shall be ignorant to it for the time being to simplify things):
  \begin{enumerate}[label=\alph*)]
    \item{\(f(x)=x\) \footnote{This function has a special name, the \textbf{identity} function. Notice how its input and output are the same.}}
    \item{\(f(x)=\frac{1}{x}\)}
    \item{\(f(x)=\sqrt{x}\)}
  \end{enumerate}

\end{problem}
\begin{problem}
  Find the root(s) of the following functions:
  \begin{enumerate}[label=\alph*)]
    \item{\(f(x)=x^2\)} 
    \item{\(f(x)=2^x-1\)}
    \item{\(f(x)=\sin(x)\)}
  \end{enumerate}
\end{problem}


\begin{problem}[1]
  A \emph{recursive} function is one where the function uses itself in its definition. For example, the Fibonacci sequence can be defined like this for any Fibonacci number \(n\): \begin{equation*}
    f(n) = \begin{cases}
      f(n-1) + f(n-2) & n>1 \\
      1 & n=1 \\
      0 & n=0 
    \end{cases}
  \end{equation*}
  \begin{enumerate}[label=\alph*)]
    \item{Compute \(f(3)\) explicitly using the definition of \(f\).}
    \item{Define the factorial of a positive integer in a similar way using recursion, define it as \(g(a)\).}
    \item{Sketch the graph of \(g(a)\).}
  \end{enumerate}

\end{problem} 

\begin{problem}[2]
  Find \(f:\mathbb{R}\setminus\{0,1\} \to  \mathbb{R}\) such that \(f(x)+f(\frac{x-1}{x})=1+x\) for any \(x\in\mathbb{R}\setminus\{0,1\}\).
\end{problem}

%     \begin{enumerate}
%       \item{
%           \begin{enumerate}
%             \item{9}
%             \item{4}
%             \item{9}
%             \item{0}
%           \end{enumerate}
%         }
%       \item{
%           \begin{enumerate}
%             \item{\(f: \mathbb{R} \to \mathbb{R}\)}
%             \item{\(f: \mathbb{R} \setminus \{0\} \to \mathbb{R} \setminus \{0\}\) (Note the function is not defined at \(x=0\))}
%             \item{\(f: \mathbb{R}^{+}_{0} \to \mathbb{R}^{+}_{0}\) (Square root not defined for negative numbers)}
%           \end{enumerate}
%       \item{
%           \begin{enumerate}
%             \item{0}
%             \item{1}
%             \item{\(n\pi, \, \forall n \in \mathbb{Z}\)}
%           \end{enumerate}
%         }
%         }
%       \item{
%           \begin{enumerate}
%       \item{
%         \(
%         f(3) = f(2) + f(1) 
%              = f(1) + f(0) + 1
%              = 2
%     \)}
%       \item{\(g(a)=\begin{cases}
%             ag(a-1) & a>0 \\
%             1 & a=0
%       \end{cases} \)}
%     \item{
%         Because our definition of \(g(a)\) is discontinuous (the function have "gaps"), we shall only plot the points where the function is defined.
%         \begin{figure}[H]
%           \centering
%       \begin{tikzpicture}
%         \begin{axis}[axis lines=left, xmin=0, xmax=5, xlabel=\(a\), ylabel=\(g(a)\)]
%           \addplot[color=blue, only marks, mark=*] coordinates{
%               (0,1)
%               (1,1)
%               (2,2)
%               (3,6)
%               (4,24)
%               (5,120)
%             }; 
%         \end{axis}
%       \end{tikzpicture}
%       \caption{Graph of \(g(a)\).}
%     \end{figure} }
% \end{enumerate}}
%     \end{enumerate}
